\chapter{DaQuVa DSL: Language Design}

\section{Introduction to DaQuVa DSL Design}
DaQuVa (Data Quality Validator) is a domain-specific language designed to facilitate the safe and efficient cleaning, validation, and analysis of large datasets within SQL databases. The language adopts a declarative programming paradigm, emphasizing "what" to do rather than "how" to do it, which aligns with the needs of data engineers and analysts who require concise, readable scripts for data quality tasks.

The DSL is tailored for post-validation data cleaning in enormous databases, where data has already undergone initial validation but may still contain inconsistencies, duplicates, or errors. DaQuVa ensures cautious operations by incorporating confidence-based scoring (e.g., dirty\_level, duplicate\_score) and explicit safety mechanisms (e.g., Here() with allowDanger) to prevent accidental data loss.

Key design principles include:
\begin{itemize}
    \item \textbf{Declarative Nature}: Users describe data transformations and validations without specifying low-level implementation details.
    \item \textbf{Safety-First}: Destructive operations require explicit thresholds and permissions.
    \item \textbf{Reusability}: Support for variables and functions to enable modular, pipeline-friendly scripts.
    \item \textbf{Extensibility}: Architecture allows integration of AI-based or rule-based tools for validation.
\end{itemize}


\section{Initial DSL Grammar}
\label{sec:grammar}

\subsection{Design Principles}

The grammar of DaQuVa is guided by three principles:

\begin{enumerate}[noitemsep]
    \item \textbf{Readability.} Statements read like short English instructions
          (\texttt{scan table."col" using tool "name"}).
    \item \textbf{Minimalism.} Only constructs required for data-quality tasks
          are included; general-purpose programming facilities are deliberately
          omitted.
    \item \textbf{Safety by default.} Destructive operations on the underlying
          database require an explicit \texttt{allowDanger} flag, making
          accidental data modification syntactically impossible.
\end{enumerate}

\subsection{Lexical Elements}
\label{sec:lexical}

The lexical layer defines the primitive tokens consumed by the parser.

\subsubsection{Keywords}

The reserved words of DaQuVa are grouped by functional category in
Table~\ref{tab:keywords}. They may not be used as identifiers.

\begin{table}[H]
\centering
\caption{Reserved keywords of DaQuVa, grouped by category.}
\label{tab:keywords}
\renewcommand{\arraystretch}{1.35}
\begin{tabular}{|p{4cm}|p{9cm}|}
\hline
\normalfont\textbf{Category} & \textbf{Keywords} \\
\hline
Database
  & \texttt{my\_db} \\
\hline
Function \& Control
  & \texttt{function},\quad \texttt{return} \\
\hline
Commands
  & \texttt{scan},\quad \texttt{filter},\quad \texttt{merge},\quad
    \texttt{output},\quad \texttt{newTable},\quad
    \texttt{findLocalDuplicates} \\
\hline
Conditions \& Logic
  & \texttt{where},\quad \texttt{AND},\quad \texttt{OR},\quad
    \texttt{columns} \\
\hline
Tool Integration
  & \texttt{using},\quad \texttt{tool} \\
\hline
Output Targets
  & \texttt{console},\quad \texttt{file},\quad \texttt{from},\quad
    \texttt{addColumns} \\
\hline
Safety Flags
  & \texttt{Here},\quad \texttt{allowDanger} \\
\hline
\end{tabular}
\end{table}

\subsubsection{Identifiers}

An identifier begins with a letter and may be followed by any combination of
letters, digits, and underscores:

\begin{lstlisting}[style=bnf]
<identifier> ::= <letter> { <letter> | <digit> | "_" }
<letter>     ::= "a" | ... | "z" | "A" | ... | "Z"
<digit>      ::= "0" | ... | "9"
\end{lstlisting}

\subsubsection{Literals}

DaQuVa supports two literal types:

\begin{lstlisting}[style=bnf]
<string> ::= '"' { <any_char_except_quote> } '"'
<number> ::= <digit> { <digit> }
\end{lstlisting}

\texttt{<number>} denotes a non-negative integer, used for thresholds and
confidence scores. String literals are delimited by double quotes and may
contain any character except an unescaped double quote.

\subsubsection{Comments}

Single-line comments begin with \texttt{\#} and extend to the end of the line.
Comments are ignored by the parser.

\subsubsection{Whitespace and Layout}

Unlike indentation-sensitive languages such as Python, DaQuVa treats
whitespace purely as a lexical separator. The parser ignores formatting
characters and does not assign semantic meaning to indentation.

The following rules apply:

\begin{itemize}[noitemsep]
    \item One or more whitespace characters may appear between tokens
          to separate them.
    \item Consecutive spaces are treated as a single separator.
    \item Tab characters and newline characters are ignored by the parser.
    \item Program structure is defined by explicit punctuation rather than
          indentation.
    \item The semicolon (\texttt{;}) is the only required statement
          terminator and separates commands.
\end{itemize}

As a result, developers may freely format DaQuVa programs for readability.
The following examples are syntactically equivalent:

\begin{lstlisting}[style=dsl]
filter potential_dirty where dirty_level > 70;
\end{lstlisting}

\begin{lstlisting}[style=dsl]
filter    potential_dirty
    where   dirty_level    >    70   ;
\end{lstlisting}

This design follows the convention of languages such as C and SQL, where
whitespace improves readability but does not influence program semantics.

\subsubsection{Operators and Punctuation}

Table~\ref{tab:entities} lists the core semantic entities of DaQuVa.

\begin{table}[H]
\centering
\caption{Core semantic entities of DaQuVa.}
\label{tab:entities}
\renewcommand{\arraystretch}{1.35}
\begin{tabular}{|p{4cm}|p{9cm}|}
\hline
\textbf{Entity} & \textbf{Description} \\
\hline

\texttt{my\_db} &
Declares the active database connection using an opaque connection string. Exactly one \texttt{my\_db} declaration is expected per program; re-declaration overwrites the current context. \\
\hline

Table &
The primary data container. Tables exist either in the connected database or as in-memory result sets produced by \texttt{scan} or \texttt{newTable}. \\
\hline

Variable &
Any command result can be bound to a variable. Variables are untyped at the syntax level but carry a \emph{result-set} value at runtime. \\
\hline

Tool &
An external validator or AI-powered analyser identified by a string name (e.g., \texttt{"fuzzy\_matcher"}, \texttt{"typo\_checker"}). The runtime resolves the name to a registered tool implementation. \\
\hline

Column / Meta-variable &
Columns such as \texttt{dirty\_level} and \texttt{duplicate\_score} are injected into result sets automatically by \texttt{scan} and \texttt{findLocalDuplicates}, respectively. They may also be added explicitly via \texttt{addColumns}. \\
\hline

Function &
A named, reusable sequence of commands parameterised by identifiers. Calling a function executes its body in a new scope; parameters shadow outer variables. \\
\hline

Threshold / Confidence &
An integer in the range $[0,100]$ representing a minimum quality score. Used in \texttt{filter} conditions and as the confidence argument to \texttt{Here()}. \\
\hline

\texttt{Here()} &
A special sink that redirects writes into the live database. Requires an optional confidence argument; when used with \texttt{merge}, it writes merged records back to the source. \\
\hline

\texttt{allowDanger} &
A boolean flag that must be present on any \texttt{merge} command that targets \texttt{Here()}. It serves as an explicit acknowledgement that the operation modifies persisted data. \\
\hline

\end{tabular}
\end{table}

\subsection{Syntactic Grammar in BNF}
\label{sec:bnf}

The complete context-free grammar of DaQuVa is given below.
Non-terminals are enclosed in angle brackets; terminals are quoted.
Curly braces denote zero-or-more repetition; square brackets denote an optional
element.

\begin{lstlisting}[style=bnf, caption={Complete BNF grammar of DaQuVa.}, label={lst:bnf}]
  <SingleSpace>     ::= " "
  
  <OneOrMoreSpaces> ::= <SingleSpace>
                      | <OneOrMoreSpaces> <SingleSpace>
  
  <opt-spaces>      ::= <OneOrMoreSpaces>
                      | <empty>
  
  <program>         ::= <opt-spaces> <statement> <opt-spaces> <program>
                      | <empty>
  
  <empty>           ::=
  
  
  <statement>       ::= <opt-spaces> <db_init> <opt-spaces>
                      | <opt-spaces> <variable_assign> <opt-spaces>
                      | <opt-spaces> <function_def> <opt-spaces>
                      | <opt-spaces> <command_statement> <opt-spaces>
  
  
  <db_init>         ::= "my_db" <string> ";"
  
  
  <variable_assign> ::= <identifier> "=" <command_statement> ";"
  
  
  <function_def>    ::= "function" <identifier>
                        "(" <opt_param_list> ")" ":"
                        <statement_block>
  
  
  <opt_param_list>  ::= <param_list>
                      | <empty>
  
  
  <param_list>      ::= <identifier> <param_list_tail>
  
  <param_list_tail> ::= "," <identifier> <param_list_tail>
                      | <empty>
  
  
  <statement_block> ::= <statement> <statement_block>
                      | <empty>
  
  
  <command_statement> ::= <scan_command>
                        | <filter_command>
                        | <find_duplicates_command>
                        | <merge_command>
                        | <output_command>
                        | <new_table_command>
                        | <function_call>
  
  
  <scan_command>    ::= "scan" <table_column> "using" "tool" <tool_name>
  
  
  <tool_name>       ::= <string>
  
  
  <table_column>    ::= <table_name> "." <column_name>
  
  <table_name>      ::= <string>
  
  <column_name>     ::= <string>
  
  
  <filter_command>  ::= "filter" <identifier> "where" <condition>
  
  
  <condition>       ::= <expression> <condition_tail>
  
  
  <condition_tail>  ::= <logical_op> <expression> <condition_tail>
                      | <empty>
  
  
  <logical_op>      ::= "AND"
                      | "OR"
  
  
  <expression>      ::= <identifier> <comparison_op> <number>
  
  
  <comparison_op>   ::= ">"
                      | "<"
                      | ">="
                      | "<="
                      | "=="
                      | "!="
  
  
  <find_duplicates_command> ::= "findLocalDuplicates" <identifier>
                                "columns" "[" <column_list> "]"
                                <opt_using_tool>
  
  
  <opt_using_tool>  ::= "using" "tool" <string>
                      | <empty>
  
  
  <column_list>     ::= <string> <column_list_tail>
  
  <column_list_tail> ::= "," <string> <column_list_tail>
                       | <empty>
  
  
  <merge_command>   ::= "merge" <identifier>
                        "->" "Here" <opt_merge_offset> <opt_allow_danger>
  
  
  <opt_merge_offset> ::= "(" <number> ")"
                       | <empty>
  
  
  <opt_allow_danger> ::= "allowDanger"
                       | <empty>
  
  
  <output_command>  ::= "output" <identifier> "->" <output_target>
  
  
  <output_target>   ::= "console"
                      | "file" <string>
                      | "Here" <opt_output_offset>
  
  
  <opt_output_offset> ::= "(" <number> ")"
                        | <empty>
  
  
  <new_table_command> ::= "newTable" <string>
                          <opt_from>
                          <opt_add_columns>
  
  
  <opt_from>        ::= "from" <identifier>
                      | <empty>
  
  
  <opt_add_columns> ::= "addColumns" "[" <column_func_list> "]"
                      | <empty>
  
  
  <column_func_list> ::= <column_func> <column_func_list_tail>
  
  <column_func_list_tail> ::= "," <column_func> <column_func_list_tail>
                            | <empty>
  
  
  <column_func>     ::= <string>
                      | <func_call>
  
  
  <func_call>       ::= <identifier> "(" <opt_arg_list> ")"
  
  
  <function_call>   ::= <identifier> "(" <opt_arg_list> ")"
  
  
  <opt_arg_list>    ::= <arg_list>
                      | <empty>
  
  
  <arg_list>        ::= <value> <arg_list_tail>
  
  <arg_list_tail>   ::= "," <value> <arg_list_tail>
                      | <empty>
  
  
  <value>           ::= <identifier>
                      | <string>
                      | <number>
  
  
  <identifier>      ::= <letter> <identifier_tail>
  
  <identifier_tail> ::= <letter> <identifier_tail>
                      | <digit> <identifier_tail>
                      | "_" <identifier_tail>
                      | <empty>
  
  
  <number>          ::= <digit> <number_tail>
  
  <number_tail>     ::= <digit> <number_tail>
                      | <empty>
  
  
  <string>          ::= '"' <string_body> '"'
  
  <string_body>     ::= <any_char_except_quote> <string_body>
                      | <empty>
  
  (* Note: <opt-spaces> and <OneOrMoreSpaces> may appear between any two tokens
     (except inside literals, identifiers, numbers and multi-character operators).
     They are shown only in the top-level rules for brevity; the rest of the grammar
     follows the same whitespace-insensitive convention. *)
  \end{lstlisting}
\subsection{Core Syntactic Structures}

\subsubsection{Program and Statements}

A DaQuVa program is a flat sequence of statements with no required entry-point.
Execution proceeds top-to-bottom. There are four statement kinds:

\begin{itemize}[noitemsep]
    \item \textbf{Database initialisation} (\texttt{my\_db}): establishes the
          active connection.
    \item \textbf{Variable assignment}: binds the result of a command to a name.
    \item \textbf{Function definition}: introduces a named, reusable routine.
    \item \textbf{Command statement}: executes an action without binding.
\end{itemize}

\subsubsection{Expressions and Conditions}

Conditions used in \texttt{filter} commands consist of one or more
\emph{expressions} joined by \texttt{AND} / \texttt{OR}. Each expression
compares a column-level attribute (e.g.\ \texttt{dirty\_level}) to an integer
threshold using a standard comparison operator. This deliberately limited
expression language prevents arbitrary logic and keeps programs auditable.

\subsubsection{Functions}

Functions accept a comma-separated list of positional parameters and contain
an indentation-delimited block of statements. Unlike general-purpose languages,
DaQuVa functions may only contain the same statement types available at the
top level; there is no nested function definition or recursion.

% ============================================================
\section{Basic DSL Semantics}
\label{sec:semantics}
% ============================================================

\subsection{Execution Model}

DaQuVa programs execute \emph{sequentially}. Each statement is fully evaluated
before the next begins. The runtime maintains:

\begin{enumerate}[noitemsep]
    \item A single \textbf{database context} set by \texttt{my\_db}.
    \item A \textbf{variable environment} mapping names to result sets.
    \item A \textbf{function table} mapping names to parameter lists and
          statement blocks.
\end{enumerate}

There is no implicit concurrency, no exception-handling syntax, and no
dynamic typing --- every variable holds a \emph{result set} (a typed, annotated
table-like structure).

\subsection{Semantic Description of Core Constructs}

Table~\ref{tab:entities} summarises the principal semantic entities of DaQuVa.

\begin{table}[H]
    \centering

\caption{Core semantic entities of DaQuVa}
\label{tab:entities}\\
\renewcommand{\arraystretch}{1.35}
\begin{tabular}{|p{3cm}|p{10cm}|}
\hline
\textbf{Entity} & \textbf{Description} \\
\hline


\texttt{my\_db} &
Declares the active database connection using an opaque connection string.
Exactly one \texttt{my\_db} declaration is expected per program; re-declaration overwrites the current context. \\
\hline

Table &
The primary data container. Tables exist either in the connected database or as in-memory result sets produced by \texttt{scan} or \texttt{newTable}. \\
\hline

Variable &
Any command result can be bound to a variable. Variables are untyped at the syntax level but carry a \emph{result-set} value at runtime. \\
\hline

Tool &
An external validator or AI-powered analyser identified by a string name (e.g. \texttt{"fuzzy\_matcher"}, \texttt{"typo\_checker"}). The runtime resolves the name to a registered tool implementation. \\
\hline

Column / Meta-variable &
Columns such as \texttt{dirty\_level} and \texttt{duplicate\_score} are injected into result sets automatically by \texttt{scan} and \texttt{findLocalDuplicates}, respectively. They may also be added explicitly via \texttt{addColumns}. \\
\hline

Function &
A named, reusable sequence of commands parameterised by identifiers. Calling a function executes its body in a new scope; parameters shadow outer variables. \\
\hline

Threshold / Confidence &
An integer in the range $[0,100]$ representing a minimum quality score. Used in \texttt{filter} conditions and as the confidence argument to \texttt{Here()}. \\
\hline

\texttt{Here()} &
A special sink that redirects writes into the live database. Requires an optional confidence argument; when used with \texttt{merge}, it writes merged records back to the source. \\
\hline

\texttt{allowDanger} &
A boolean flag that must be present on any \texttt{merge} command that targets \texttt{Here()}. It serves as an explicit acknowledgement that the operation modifies persisted data. \\
\hline
\end{tabular}
\end{table}

\subsection{Command Semantics}

\subsubsection{\texttt{scan} --- Column Validation}

\begin{lstlisting}[style=bnf]
<scan_command> ::= "scan" <table_column> "using" "tool" <string>
\end{lstlisting}

\textbf{Meaning.} \texttt{scan} passes every value in the named column of the
named table to the specified tool. The tool assigns each row a
\texttt{dirty\_level} score in $[0, 100]$, where 100 represents maximum
dirtiness. The result is an annotated copy of the rows with the added
meta-column.

\textbf{Execution.}
\begin{enumerate}[noitemsep]
    \item Resolve the table name against the active database context.
    \item Resolve the tool name against the registered tool registry.
    \item Invoke the tool for each row value in the column.
    \item Return a result set containing all original columns plus
          \texttt{dirty\_level}.
\end{enumerate}

\textbf{Errors.} An error is raised if the table or column does not exist,
if the tool name cannot be resolved, or if the database context has not been
initialised.

\subsubsection{\texttt{filter} --- Threshold Filtering}

\begin{lstlisting}[style=bnf]
<filter_command> ::= "filter" <identifier> "where" <condition>
\end{lstlisting}

\textbf{Meaning.} Produces a subset of the input result set containing only
rows that satisfy the given condition.

\textbf{Execution.} Each row is evaluated against the condition; rows for which
the condition is true are retained. Conditions may combine multiple predicates
with \texttt{AND} / \texttt{OR}; evaluation follows standard boolean semantics
with left-to-right associativity.

\textbf{Errors.} An error is raised if the referenced variable is undefined,
if a column referenced in a condition does not exist in the result set, or if
operand types are incompatible.

\subsubsection{\texttt{findLocalDuplicates} --- Duplicate Detection}

\begin{lstlisting}[style=bnf]
<find_duplicates_command>
    ::= "findLocalDuplicates" <identifier>
            "columns" "[" <column_list> "]"
            [ "using" "tool" <string> ]
\end{lstlisting}

\textbf{Meaning.} Identifies rows within the input result set that are
potential duplicates according to the specified key columns. An optional
AI-powered tool (e.g.\ a fuzzy matcher) provides similarity scoring.

\textbf{Execution.} Rows are grouped by the specified columns; within each
group, pairs are scored for similarity. The result set contains all candidate
duplicate rows annotated with a \texttt{duplicate\_score} in $[0, 100]$.

\textbf{Errors.} An error is raised if the input variable is undefined or if
any listed column does not exist in the result set.

\subsubsection{\texttt{merge} --- Duplicate Resolution}

\begin{lstlisting}[style=bnf]
<merge_command> ::= "merge" <identifier>
                        "->" "Here" [ "(" <number> ")" ]
                        [ "allowDanger" ]
\end{lstlisting}

\textbf{Meaning.} Merges duplicate records in the input result set and
writes the resolved records back to the live database through \texttt{Here()}.
The optional numeric argument specifies a minimum confidence threshold;
only duplicate pairs with \texttt{duplicate\_score} $\geq$ threshold are
merged.

\textbf{Execution.}
\begin{enumerate}[noitemsep]
    \item Filter candidates by the confidence threshold (default: 0).
    \item For each qualifying pair, apply a deterministic merge strategy
          (e.g.\ keep the most complete record).
    \item Write merged records to the database; delete superseded records.
\end{enumerate}

\textbf{Safety constraint.} The \texttt{allowDanger} flag \emph{must} be
present. Its absence is a compile-time error. This design makes destructive
database modifications opt-in and visible in code review.

\subsubsection{\texttt{output} --- Result Export}

\begin{lstlisting}[style=bnf]
<output_command> ::= "output" <identifier> "->" <output_target>
<output_target>  ::= "console" | "file" <string> | "Here" [ "(" <number> ")" ]
\end{lstlisting}

\textbf{Meaning.} Writes a result set to one of three sinks: the console
(standard output), a named file, or the live database via \texttt{Here()}.

\textbf{Execution.}
\begin{itemize}[noitemsep]
    \item \texttt{console} --- serialises rows as human-readable text.
    \item \texttt{file "<path>"} --- writes rows to the specified file path in
          a delimited format (e.g.\ CSV).
    \item \texttt{Here()} --- inserts or upserts rows into the database.
\end{itemize}

\textbf{Errors.} An error is raised if the variable is undefined, if the
file path is not writable, or if the database context is missing for
\texttt{Here()}.

\subsubsection{\texttt{newTable} --- Table Construction}

\begin{lstlisting}[style=bnf]
<new_table_command> ::= "newTable" <string>
                            [ "from" <identifier> ]
                            [ "addColumns" "[" <column_func_list> "]" ]
\end{lstlisting}

\textbf{Meaning.} Creates a new named table, optionally seeded from an
existing result set and optionally extended with computed columns.

\textbf{Execution.} If \texttt{from} is provided, the new table is initialised
with the rows of the referenced result set. Each entry in \texttt{addColumns}
is either a plain column name (copied from the source) or a function call
whose result becomes a new column value (e.g.\ \texttt{length("name")}).

\subsubsection{\texttt{function} and Function Calls}

\begin{lstlisting}[style=bnf]
<function_def>  ::= "function" <identifier> "(" [ <param_list> ] ")" ":"
                        <statement_block>
<function_call> ::= <identifier> "(" [ <arg_list> ] ")"
\end{lstlisting}

\textbf{Meaning.} Functions encapsulate reusable command sequences. They are
\emph{first-defined} (must be defined before use in a top-to-bottom
interpreter) and execute in their own variable scope. Parameters are
positional; the number of arguments at the call site must match the parameter
list.

\textbf{Errors.} Arity mismatch, calling an undefined function, or a
\texttt{return} statement referencing an undefined variable all raise errors.

\subsection{Error Catalogue}

Table~\ref{tab:errors} catalogues the principal errors and their triggers.

\begin{table}[H]
\centering
\caption{Error catalogue for DaQuVa.}
\label{tab:errors}
\renewcommand{\arraystretch}{1.35}
\begin{tabular}{|p{4cm}|p{10cm}|}
\hline
\normalfont\textbf{Error} & \textbf{Condition} \\
\hline
UndefinedVariable & An identifier is used before assignment. \\
\hline
UndefinedTable    & A table name cannot be resolved in the database context. \\
\hline
UndefinedColumn   & A column referenced in a condition or column list does
                    not exist in the result set. \\
\hline
UnresolvedTool    & A tool string cannot be matched to a registered tool. \\
\hline
ArityMismatch     & A function call provides more or fewer arguments than
                    parameters declared. \\
\hline
MissingContext    & A database-dependent command is issued before
                    \texttt{my\_db}. \\
\hline
MissingAllowDanger & A \texttt{merge -> Here()} command lacks the
                    \texttt{allowDanger} flag. \\
\hline
InvalidThreshold  & A numeric argument to \texttt{Here()} or a filter
                    expression falls outside $[0, 100]$. \\
\hline
IOError           & The file path in an \texttt{output -> file} command
                    is not accessible. \\
\hline
\end{tabular}
\end{table}

% ============================================================
\section{Illustrative Example Program}
\label{sec:example}
% ============================================================

The program below demonstrates all major language constructs. It connects to a
database, defines a reusable email-cleaning function, applies it, exports
results, detects and merges duplicates, and creates a summary table for further
analysis.

\begin{lstlisting}[style=dsl, caption={Annotated DaQuVa example program.}, label={lst:example}]
# === 1. Database connection ===
my_db "some_connection_string_goes_here";

# === 2. Reusable function: clean emails in any table ===
function cleanEmails(table_name, threshold):
    # Scan the email column with an AI validator
    potential_dirty = scan table_name."email" using tool "misa_ai_validator";

    # Keep only rows whose dirty_level exceeds the threshold
    high_risk = filter potential_dirty where dirty_level > threshold;

    return high_risk;

# === 3. Apply the function to the "users" table ===
high_risk_emails = cleanEmails("users", 70);

# === 4. Export results ===
output high_risk_emails -> file "high_risk_emails.csv";
output high_risk_emails -> console;

# === 5. Detect near-duplicate records among high-risk rows ===
duplicates = findLocalDuplicates high_risk_emails
             columns ["name", "dob"]
             using tool "fuzzy_matcher";

# === 6. Merge duplicates back into the database (confidence >= 90%) ===
merge duplicates -> Here(90) allowDanger;

# === 7. Create a summary table for downstream analysis ===
analysis_table = newTable "high_risk_analysis"
                 from high_risk_emails
                 addColumns [duplicate_score, length("name")];

# === 8. Additional validation pass on the summary table ===
cli_result = scan analysis_table."email" using tool "typo_checker";
output cli_result -> console;
\end{lstlisting}

The program illustrates several semantic properties:

\begin{itemize}
    \item The function \texttt{cleanEmails} abstracts a two-step
          scan-then-filter pipeline, making the threshold a tuneable parameter.
    \item The \texttt{output} command is used twice on the same variable,
          demonstrating that result sets are non-destructively consumed.
    \item The \texttt{merge} command requires \texttt{allowDanger}
          because it writes to the live database, making the risky operation
          immediately visible to reviewers.
    \item \texttt{newTable} constructs a derived table that mixes copied data
          with computed columns (\texttt{duplicate\_score}, \texttt{length("name")}),
          enabling further analysis without modifying the source table.
\end{itemize}

% ============================================================
\section{Derivation Trees for Control Structures}
% ============================================================

Below are the derivation trees for the main control structures of DaQuVa.

% ----------------------------------------------------------
\subsection{Variable Definition}
% ----------------------------------------------------------

\textbf{Example:} \texttt{high\_risk\_emails = cleanEmails("users", 70);}

\begin{center}
\begin{adjustbox}{max width=\textwidth}
\begin{forest} daqtree
[statement
  [variable\_assign
    [identifier [high\_risk\_emails]]
    [=]
    [function\_call
      [cleanEmails]
      [arg\_list
        ["users"]
        [,]
        [70]
      ]
    ]
    [;]
  ]
]
\end{forest}
\end{adjustbox}
\end{center}

% ----------------------------------------------------------
\subsection{scan Command}
% ----------------------------------------------------------

\textbf{Example:} \texttt{scan table\_name."email" using tool "misa\_ai\_validator";}

\begin{center}
\begin{adjustbox}{max width=\textwidth}
\begin{forest} daqtree
[statement
  [scan\_command
    [scan]
    [table\_column
      [table\_name]
      [.]
      ["email"]
    ]
    [using]
    [tool]
    ["misa\_ai\_validator"]
    [;]
  ]
]
\end{forest}
\end{adjustbox}
\end{center}

% ----------------------------------------------------------
\subsection{filter Command}
% ----------------------------------------------------------

\textbf{Example:} \texttt{filter potential\_dirty where dirty\_level > 70;}

\begin{center}
\begin{adjustbox}{max width=\textwidth}
\begin{forest} daqtree
[statement
  [filter\_command
    [filter]
    [potential\_dirty]
    [where]
    [condition
      [dirty\_level]
      [>]
      [70]
    ]
    [;]
  ]
]
\end{forest}
\end{adjustbox}
\end{center}

% ----------------------------------------------------------
\subsection{Function Definition}
% ----------------------------------------------------------

\textbf{Example:} \texttt{function cleanEmails(table\_name, threshold):}

\begin{center}
\begin{adjustbox}{max width=\textwidth}
\begin{forest} daqtree
[function\_def
  [function]
  [cleanEmails]
  [param\_list
    [table\_name]
    [,]
    [threshold]
  ]
  [:]
  [statement\_block]
]
\end{forest}
\end{adjustbox}
\end{center}

% ----------------------------------------------------------
\subsection{merge Command}
% ----------------------------------------------------------

\textbf{Example:} \texttt{merge duplicates -> Here(90) allowDanger;}

\begin{center}
\begin{adjustbox}{max width=\textwidth}
\begin{forest} daqtree
[statement
  [merge\_command
    [merge]
    [duplicates]
    [->]
    [Here]
    [(90)]
    [allowDanger]
    [;]
  ]
]
\end{forest}
\end{adjustbox}
\end{center}

% ----------------------------------------------------------
\subsection{output Command}
% ----------------------------------------------------------

\textbf{Example:} \texttt{output high\_risk\_emails -> file "high\_risk\_emails.csv";}

\begin{center}
\begin{adjustbox}{max width=\textwidth}
\begin{forest} daqtree
[statement
  [output\_command
    [output]
    [high\_risk\_emails]
    [->]
    [file]
    ["high\_risk\_emails.csv"]
    [;]
  ]
]
\end{forest}
\end{adjustbox}
\end{center}

% ----------------------------------------------------------
\subsection{findLocalDuplicates Command}
% ----------------------------------------------------------

\textbf{Example:} \texttt{findLocalDuplicates high\_risk\_emails columns ["name","dob"];}

\begin{center}
\begin{adjustbox}{max width=\textwidth}
\begin{forest} daqtree
[statement
  [find\_duplicates\_command
    [findLocalDuplicates]
    [high\_risk\_emails]
    [columns]
    [column\_list
      ["name"]
      [,]
      ["dob"]
    ]
    [;]
  ]
]
\end{forest}
\end{adjustbox}
\end{center}

% ----------------------------------------------------------
\subsection{newTable Command}
% ----------------------------------------------------------

\textbf{Example:} \texttt{newTable "high\_risk\_analysis" from high\_risk\_emails addColumns [duplicate\_score];}

\begin{center}
\begin{adjustbox}{max width=\textwidth}
\begin{forest} daqtree
[statement
  [new\_table\_command
    [newTable]
    ["high\_risk\_analysis"]
    [from]
    [high\_risk\_emails]
    [addColumns]
    [duplicate\_score]
    [;]
  ]
]
\end{forest}
\end{adjustbox}
\end{center}


% ============================================================
\section{Summary}
% ============================================================

This report has presented the complete formal grammar and semantic
specification of the \textbf{Data Quality Validator} DSL. The BNF grammar
in Section~\ref{sec:grammar} defines an unambiguous, minimal syntax
covering database initialisation, column scanning, result filtering,
duplicate detection and merging, result export, and table construction.
The semantic specification in Section~\ref{sec:semantics} assigns a precise
meaning to each construct, documents execution rules, and enumerates the
error conditions the runtime must detect and report.

The language's design deliberately constrains expressiveness in favour of
safety and auditability: destructive operations require explicit opt-in flags,
and the expression language is limited to simple threshold comparisons,
preventing opaque side effects. These properties make DaQuVa programs
straightforward to review, test, and extend.